\section{Concluding Discussion}
\label{sec:conclusion}
\label{sec:recs}
\label{sec:limits}

Our study evaluates self-recognition in nine LLM judges across three sets of
texts and three question formats. The usual test, a self-naming rate above
chance, flags self-recognition in \SVPStd{} cases, and only \SVPProper{} survive
a peer baseline and a false-alarm rate. Where self-recognition does exist, it
is large and invisible to anyone who looks at only one format or one set of
texts: Claude and GPT beat their peers by 33 and 28 points on chat-app essays
but not on 600 API texts, and the question format alone produces self-denial in
some judges and over-attribution in others. The information a judge would need
to recognise itself is already in the text and in its own probabilities, yet
the judges' answers do not follow it. Our contribution is not to declare
anonymisation safe or broken, but to provide a test that separates the two and
can be applied to any judging pipeline.

\paragraph{Implications}
Self-recognition is not a fixed trait of a model that can be certified once. It
depends on the judge, the texts and the question, so a judge that passes an
audit on one benchmark may leak on another, and a model update may change the
answer. Much of the work on self-preference assumes that judges recognise
their own outputs; our results suggest that, for most judges, the preference
must come from elsewhere, consistent with \citet{davidson2024anonymized} and
\citet{bai2025selfrecognition}. For the judges that do recognise themselves,
however, anonymisation offers no protection, and the usual test cannot tell
these judges from the many that merely appear to.

\paragraph{Recommendations}
We recommend four practices. (1)~Report the self-naming rate together with the
false-alarm rate and the peer baseline, which come from the same experiment,
and claim self-recognition only if a judge both discriminates and beats its
peers. (2)~Prefer a lineup to the one-text format, in which judges collapse onto
one or two names; use a yes/no question only if a judge's answers vary.
(3)~Measure each judge on the kind of text it will judge, and repeat the audit
when the judge or the texts change. (4)~Do not infer the mechanism from the
judges' stated reasons, which are the same for right and wrong self-attributions.
Our released code implements this audit.

\paragraph{Future directions}
The gap between what the texts reveal and what judges use suggests that
self-recognition could grow quickly if judges learned to use stylistic or
likelihood cues, through training, few-shot examples \citep{dietrich2026splits}
or access to their own probabilities. Tracking self-recognition across model
versions is therefore as important as measuring it once. Our test extends
directly to other judging tasks, such as pairwise preference and safety
monitoring, and to longer and multilingual texts. Interpretability work on the
representations behind self-attribution \citep{ackerman2025inspection} could
show why some judges discriminate on some genres and not others.

\paragraph{Limitations}
The chat-app essays contain only 38 texts per judge, so their intervals are
wide; the API texts have 120 per judge. The open-weight judges are small (4--8
billion parameters, run in 4-bit) and much weaker than frontier models. No
frontier provider returned token probabilities, so the frontier yes/no results
are hard answers, and the likelihood, quality and answer-first analyses cover
the open-weight judges only. All texts are in English and every prompt fixes the
content of the answer. The chat-app essays were collected in consumer apps whose
hidden system prompts we cannot see. The coding of the judges' reasons was
checked against two LLM annotators, both instances of Claude, which is also a
judge. Finally, every result describes particular model versions at particular
dates in 2026; providers update models behind stable names, and the DeepSeek
endpoint could not be pinned.

\section*{Acknowledgments}
We thank Ryan Erlikhman for his work on an earlier version of this study.

\section*{Reproducibility Statement}
Section~\ref{sec:setup} describes how we built the texts and called the models,
Appendix~\ref{app:data} gives the collection details, and
Appendix~\ref{app:instructions} reproduces every instruction given to a judge.
All texts, prompts, judgments with their reasons, archived API responses,
open-weight likelihoods and quality ratings, and the analysis code are available
at \CodeURL{}. Every number in the paper is generated from those files, and none
of the analyses requires an API key.

\section*{Ethics Statement}
The study involves no human subjects or personal data: every text was written
by a language model in response to our prompts. The frontier texts were
collected through the providers' apps and APIs under their terms of use, and the
open-weight models were run under their licences. We name the judges that
recognise their own text so that practitioners can account for them; the
results describe specific model versions and should not be read as permanent
properties of a vendor's models.
